% Options for packages loaded elsewhere
\PassOptionsToPackage{unicode}{hyperref}
\PassOptionsToPackage{hyphens}{url}
%
\documentclass[
  11pt,
]{article}
\usepackage{lmodern}
\usepackage{amssymb,amsmath}
\usepackage{ifxetex,ifluatex}
\ifnum 0\ifxetex 1\fi\ifluatex 1\fi=0 % if pdftex
  \usepackage[T1]{fontenc}
  \usepackage[utf8]{inputenc}
  \usepackage{textcomp} % provide euro and other symbols
\else % if luatex or xetex
  \usepackage{unicode-math}
  \defaultfontfeatures{Scale=MatchLowercase}
  \defaultfontfeatures[\rmfamily]{Ligatures=TeX,Scale=1}
\fi
% Use upquote if available, for straight quotes in verbatim environments
\IfFileExists{upquote.sty}{\usepackage{upquote}}{}
\IfFileExists{microtype.sty}{% use microtype if available
  \usepackage[]{microtype}
  \UseMicrotypeSet[protrusion]{basicmath} % disable protrusion for tt fonts
}{}
\makeatletter
\@ifundefined{KOMAClassName}{% if non-KOMA class
  \IfFileExists{parskip.sty}{%
    \usepackage{parskip}
  }{% else
    \setlength{\parindent}{0pt}
    \setlength{\parskip}{6pt plus 2pt minus 1pt}}
}{% if KOMA class
  \KOMAoptions{parskip=half}}
\makeatother
\usepackage{xcolor}
\IfFileExists{xurl.sty}{\usepackage{xurl}}{} % add URL line breaks if available
\IfFileExists{bookmark.sty}{\usepackage{bookmark}}{\usepackage{hyperref}}
\hypersetup{
  hidelinks,
  pdfcreator={LaTeX via pandoc}}
\urlstyle{same} % disable monospaced font for URLs
\usepackage[margin=1in]{geometry}
\setlength{\emergencystretch}{3em} % prevent overfull lines
\providecommand{\tightlist}{%
  \setlength{\itemsep}{0pt}\setlength{\parskip}{0pt}}
\setcounter{secnumdepth}{-\maxdimen} % remove section numbering
\usepackage{mathrsfs}
\usepackage{esint}
\setmathfont{Asana Math}
\setlength{\emergencystretch}{3em}

\author{}
\date{23 September 2026}

\begin{document}

\hypertarget{integral-einstein-rigidity-from-polarized-monopole-classes}{%
\section{Integral Einstein rigidity from polarized monopole
classes}\label{integral-einstein-rigidity-from-polarized-monopole-classes}}

\hypertarget{abstract}{%
\subsection{Abstract}\label{abstract}}

We prove a scale-invariant, integral obstruction to Einstein metrics on
closed oriented four-manifolds with suitable polarized Seiberg--Witten
classes. The analytic statement uses actual solutions for every positive
weight, but requires no Sobolev, injectivity-radius, diameter, volume,
pointwise-curvature or curvature-derivative bound. Its proof retains the
nonnegative errors in the monopole inequalities, bounds the spinor by a
truncated weight, constructs a globally smooth polynomial Weyl model,
and tests the contracted Bianchi identity against the divergence of that
model. A bundle Hessian estimate proved below closes this self-test
without a geometric compactness argument.

The resulting universal constant gives an integral Einstein gap for the
connected sums occurring in LeBrun's Einstein obstruction. The class may
depend on the metric, including its chamber when the positive Betti
number is one; the same selected class must work for every weight in
that metric's conformal class. Positive and negative Einstein constants
and the separate Ricci-flat case are treated explicitly. The
positive-budget application strengthens the numerical obstruction in
LeBrun's Theorem 3.3 and Corollary 3.4. Negative topological budgets are
handled directly by Chern--Gauss--Bonnet; the zero-budget, zero-surgery
corner is not an Einstein nonexistence statement.

This manuscript makes the inherited integral-rigidity argument
self-contained analytically. It does not assert priority for either that
argument or the classical gauge-theoretic inputs.

\hypertarget{conventions-and-statements}{%
\subsection{1. Conventions and
statements}\label{conventions-and-statements}}

Let \(M\) be a closed, connected, oriented smooth four-manifold, and let

\[
 V=\operatorname{Vol}_g(M),\qquad T=2\chi(M)+3\tau(M),\qquad
 P=\int_M|\operatorname{Ric}^{\circ}|^2.
\]

Integrals and norms use \(g\), unless stated otherwise. Ricci is a
covariant two-tensor; \(W^+\) is the trace-free symmetric curvature
endomorphism of \(\Lambda^+\), with its operator Hilbert--Schmidt norm.
The Laplacian is nonnegative. Write \(x_-:=\max\{-x,0\}\) and

\[
 4\pi^2T=\int_M\left(\frac{s^2}{24}+2|W^+|^2
                          -\frac{|\operatorname{Ric}^{\circ}|^2}{2}\right).
 \tag{1.1}
\]

For a Spin\(^c\) structure with determinant line \(L\), put

\[
 a=c_1(L)\in H^2(M;\mathbb R),\qquad
 a^+=\operatorname{pr}_{H_g^+}a,\qquad p=(a^+)^2.
\]

All cohomological squares mean intersection squares. The quadratic
spinor map is normalized by

\[
 |\sigma(\Phi)|^2=|\Phi|^4/8.
\]

\textbf{Definition 1.1 (the conformal monopole hypothesis).} A pair
\((g,a)\) satisfies this hypothesis if \(p>0\) and one fixed Spin\(^c\)
structure with first Chern class \(a\) has an unperturbed
Seiberg--Witten solution for every metric conformal to \(g\).
Equivalently for our purposes, for every smooth \(f>0\) there is a
solution of

\[
 D_A\Phi=0,\qquad -iF_A^+=f\sigma(\Phi).                 \tag{1.2}
\]

The implication from conformal solvability to (1.2) is proved in §2. The
class is not allowed to change as \(f\) changes. A different class may
be selected for a different initial metric.

Define the scale-invariant integral Einstein error by

\[
 \mathcal E(g):=\inf_{\lambda\ne0}
  \frac{\int_M|\operatorname{Ric}_g-\lambda g|^2}
       {\lambda^2V}.                                      \tag{1.3}
\]

\textbf{Theorem 1.2 (integral rigidity).} There is an absolute constant

\[
                       \varepsilon_0>0                   \tag{1.4}
\]

with the following property. Suppose that for a metric \(g\) there is a
class \(a\) satisfying Definition 1.1 and

\[
                     (a^+)^2\ge A>0,\qquad a^2\ge T.
 \tag{1.5}
\]

Then

\[
 \boxed{\displaystyle
 \mathcal E(g)+\left(\frac TA-\frac23\right)_+
                         \ge\varepsilon_0.}              \tag{1.6}
\]

In particular, the conclusion holds for every metric if such a class can
be selected metricwise, with the same number \(A\). No regularity of the
class selection is needed. The constant is existential and universal;
this paper does not assign it a numerical value.

The main analytic ingredient is a calibration gap. Let

\[
 \begin{split}
 \mathcal O&=\operatorname{Orb}\operatorname{diag}
                    (-1/6,1/12,1/12)\subset S^2_0(\mathbb R^3),\\
 F&=\operatorname{dist}_{\mathrm{HS}}(W^+,\mathcal O)^2,\qquad
 E=\operatorname{Ric}+g,\qquad e=s+4,\\
 \alpha&=\sqrt{2/3}|W^+|,\qquad q=-s+3\alpha,\\
 \mathcal C&=\int_M(|E|^2+F),\qquad
 \Delta_+=\int_Mq_+^2-72\pi^2p,\\
 \kappa&=(\mathcal C+\Delta_+)/V,\qquad
 B_a=(-a^2)_+/V.
 \end{split}                                              \tag{1.7}
\]

Here \(\mathcal O\) is interpreted in each oriented orthonormal frame of
\(\Lambda^+\); the distance is frame independent.

\textbf{Theorem 1.3 (analytic calibration gap).} Under Definition 1.1,

\[
 \Delta_+\ge0,\qquad (1+B_a)\kappa\ge c_{\mathrm{cal}}>0, \tag{1.8}
\]

where \(c_{\mathrm{cal}}\) is absolute.

We prove this theorem first. Throughout the analytic proof, \(C\)
denotes an absolute positive constant whose value can change from line
to line. Explicit intermediate constants are retained where they clarify
the topological transfer, not to advertise an optimized final constant.

\hypertarget{actual-weighted-solutions-and-the-retained-monopole-identities}{%
\subsection{2. Actual weighted solutions and the retained monopole
identities}\label{actual-weighted-solutions-and-the-retained-monopole-identities}}

\hypertarget{why-the-weighted-equation-has-solutions}{%
\subsubsection{2.1. Why the weighted equation has
solutions}\label{why-the-weighted-equation-has-solutions}}

Let \(\widehat g=e^{2u}g=f^{-2}g\), so \(u=-\log f\). Under the
fiberwise unitary spinor identification \(\beta\), the conformal
formulas are

\[
 \begin{split}
 D_A^{\widehat g}(e^{-3u/2}\beta\Phi)
        &=e^{-5u/2}\beta D_A^g\Phi,\\
 \sigma_{\widehat g}(\beta v)&=e^{2u}\sigma_g(v)
                   \quad\hbox{as covariant two-forms}.
 \end{split}                                              \tag{2.1}
\]

For an ordinary solution \((A,\widehat\Phi)\) for \(\widehat g\), set

\[
                 \Phi=f^{-3/2}\beta^{-1}\widehat\Phi.
\]

Then \(D_A^g\Phi=0\), and

\[
 \sigma_g(\Phi)=f^{-1}\sigma_{\widehat g}(\widehat\Phi).
\]

The Hodge star on two-forms is conformally invariant in dimension four,
so the ordinary curvature equation becomes (1.2). In particular
\(H_g^+\) does not change during this operation. This is the conformal
weighted monopole construction of
\href{https://arxiv.org/pdf/math/0611450}{LeBrun, Lemma 3.7}; the
displayed formulas specify the normalization used here.

Set

\[
             \rho=|\Phi|^2,\qquad \psi=2\sqrt2\sigma(\Phi).
 \tag{2.2}
\]

Thus \(\psi\) is a smooth self-dual two-form, \(|\psi|=\rho\), and
\(\rho\) itself is smooth, including at spinor zeros. The harmonic
projection of the determinant curvature gives

\[
 \int_M f^2\rho^2=8\|F_A^+\|_2^2\ge32\pi^2p.           \tag{2.3}
\]

Only the squared norm is used, so the equivalent opposite sign
convention for the Chern representative does not change this formula.

\hypertarget{nonnegative-errors}{%
\subsubsection{2.2. Nonnegative errors}\label{nonnegative-errors}}

The spinor Weitzenböck equation and the self-dual Hodge formula are

\[
 4\nabla_A^*\nabla_A\Phi+(s+f\rho)\Phi=0,                \tag{2.4}
\]

\[
 \|d\psi\|_2^2+\|d^*\psi\|_2^2
   =\int_M\left(|\nabla\psi|^2+\frac{s}{3}\rho^2
                                  -2\langle W^+\psi,\psi\rangle\right).
 \tag{2.5}
\]

The differential of the Hopf quadratic map in (2.2) satisfies

\[
 |d\psi_z(v)|^2
       =4|z|^2|v|^2-4\bigl(\operatorname{Im}\langle z,v\rangle\bigr)^2
       \le4|z|^2|v|^2.                                   \tag{2.6}
\]

The same identity applies covariantly because the map is equivariant and
parallel. Pairing (2.4) with \(\rho\Phi\) and integrating yields

\[
 2\int_M|d\rho|^2+4\int_M\rho|\nabla_A\Phi|^2
                  +\int_Ms\rho^2+\int_M f\rho^3=0.       \tag{2.7}
\]

Let \(w\) be the smallest eigenvalue of \(W^+\). For a trace-free
symmetric three-by-three matrix, \(w\ge-\alpha\). Define

\[
 \begin{split}
 D_{\mathrm{amp}}&=2\int_M|d\rho|^2,\\
 D_{\mathrm{quad}}&=\int_M(4\rho|\nabla_A\Phi|^2-|\nabla\psi|^2),\\
 D_H&=\|d\psi\|_2^2+\|d^*\psi\|_2^2=2\|d\psi\|_2^2,\\
 D_{\mathrm{eig}}&=2\int_M
           (\langle W^+\psi,\psi\rangle+\alpha\rho^2),\\
 D_0&=D_{\mathrm{amp}}+D_{\mathrm{quad}}+D_H+D_{\mathrm{eig}}.
 \end{split}                                              \tag{2.8}
\]

Each term is nonnegative. Equations (2.5)--(2.7) give the exact identity

\[
 D_0=\int_M\left(\frac23q\rho^2-f\rho^3\right).
 \tag{2.9}
\]

For \(v=3\rho/2\), this is

\[
                  \int_M(qv^2-fv^3)=\frac{27}{8}D_0.
 \tag{2.10}
\]

\hypertarget{signed-spectral-and-scalar-estimates}{%
\subsubsection{2.3. Signed spectral and scalar
estimates}\label{signed-spectral-and-scalar-estimates}}

Put

\[
                  b=2s/3+2w,\qquad h=b_-.
\]

Using \(w\), rather than \(-\alpha\), in (2.5)--(2.7) gives

\[
                       \int_M f\rho^3\le\int_M h\rho^2.
 \tag{2.11}
\]

Choose positive smooth functions \(h_j\ge h\) converging uniformly to
\(h\), and use the actual solutions with \(f_j=\sqrt{h_j}\). Set

\[
 C_j=\int_M f_j^2\rho^2,\quad Z_j=\int_M f_j\rho^3,
 \quad Q_j=\int_M f_j^4.
\]

Equation (2.11) and Hölder imply

\[
 32\pi^2p\le C_j\le Q_j^{1/3}Z_j^{2/3}
                            \le Q_j^{1/3}C_j^{2/3}.
\]

Consequently \(32\pi^2p\le Q_j\), and uniform approximation proves

\[
 \int_M(2s/3+2w)_-^2\ge32\pi^2p,\qquad
 \int_M q_+^2\ge72\pi^2p.                               \tag{2.12}
\]

The second follows from \(b_-\le(2/3)q_+\). This argument passes to a
limit of inequalities, not to a limiting spinor or connection.

For \(f=1\), pairing (2.4) with \(\Phi\) gives

\[
 \int_M\rho^2\le\int_Ms_-\rho
        \le\|s_-\|_2\|\rho\|_2.
\]

Together with (2.3), this proves

\[
                        \int_Ms_-^2\ge32\pi^2p.          \tag{2.13}
\]

In particular \(p>0\) excludes a scalar-flat metric. These are the
classical monopole curvature estimates in the normalization needed here;
the retained terms in (2.8) will supply their quantitative use.

\hypertarget{a-bounded-spinor-with-small-retained-errors}{%
\subsection{3. A bounded spinor with small retained
errors}\label{a-bounded-spinor-with-small-retained-errors}}

The weights used in §2.3 need not bound the spinor uniformly. For the
calibration argument we instead set

\[
       Q=\max\{q,3\},\qquad
       f_0=\max\{\sqrt Q,s_-/4\}.                        \tag{3.1}
\]

Choose smooth majorants \(f_j\ge f_0\) with \(f_j\to f_0\) uniformly,
and choose an actual solution for each \(f_j\). At a positive maximum of
\(\rho\), (2.4) implies \(s+f_j\rho\le0\). Hence

\[
                           0\le\rho\le4,\qquad 0\le v\le6.
 \tag{3.2}
\]

All errors \(o_j(1)\) below tend to zero for this fixed smooth metric
and fixed class. They require only uniform approximation of \(f_0\), not
bounds on derivatives of \(f_j\). No assertion of convergence of the
solutions is made.

The two truncations cost only scalar error:

\[
 0\le Q^2-q_+^2\le9e^2,\qquad
 0\le f_0^2-Q\le\frac9{16}e^2.                          \tag{3.3}
\]

Indeed, the first difference vanishes unless \(q<3\), which implies
\(s>-3\) and \(e>1\). For the second, a nonzero difference implies
\(s_->4\sqrt3>6\), and then \(s_-\le3|s+4|\).

Define ideal moments using \(Q\), not the full truncated weight:

\[
 C_*:=\int_M Qv^2,\qquad Z_*:=\int_M\sqrt Qv^3,
 \qquad E_Q:=\int_M Q^2,\qquad P_*:=C_*-Z_*.
 \tag{3.4}
\]

Equation (2.10), \(Q\ge q\), and \(f_j\ge\sqrt Q\) give

\[
 P_*\ge\frac{27}{8}D_0\ge0.                            \tag{3.5}
\]

The actual curvature moment is

\[
 C_{\mathrm{act}}:=\int_M f_j^2v^2
                      =18\|F_A^+\|_2^2\ge72\pi^2p.
\]

By (3.2)--(3.3),

\[
 0\le C_{\mathrm{act}}-C_*
                     \le\frac{81}{4}\int_M e^2+o_j(1).
\]

Hölder gives \(C_*\le E_Q^{1/3}Z_*^{2/3}\). Since \(Z_*\le C_*\), it
follows that \(C_*\le E_Q\). Thus, with

\[
                        S_*:=\Delta_++\frac{117}{4}\int_Me^2,
\]

we have

\[
                     0\le E_Q-C_*\le S_*+o_j(1).         \tag{3.6}
\]

For completeness, two elementary moment inequalities recover all the
errors needed below. Write

\[
 \theta=(Z_*/E_Q)^{1/3},\qquad
 H_*=E_Q^{1/3}Z_*^{2/3}-C_*.
\]

Here \(0<\theta\le1\), \(H_*\ge0\), and direct algebra gives

\[
 E_Q-C_*=H_*+\frac{1+\theta}{\theta^2}(P_*+H_*)\ge2P_*.
 \tag{3.7}
\]

For the probability measure \(d\nu=Q^2d\mu/E_Q\), let \(r=v/\sqrt Q\),
\(c=\int r^2d\nu\), \(z=\int r^3d\nu\), and \(m=\int rd\nu\). We have
\(z\le c\le z^{2/3}\), hence \(c\le1\). Also \(c^2\le mz\le mc\), so
\(m\ge c\). Therefore

\[
 \int_M Q(v-\sqrt Q)^2=E_Q(c-2m+1)\le E_Q-C_*.
 \tag{3.8}
\]

Equations (3.5)--(3.8) prove

\[
 \begin{split}
 D_0&\le\frac4{27}S_*+o_j(1),\\
 \int_M|d\rho|^2,\ \int_M|d\psi|^2
     &\le\frac2{27}S_*+o_j(1),\\
 \int_MQ(v-\sqrt Q)^2&\le S_*+o_j(1).
 \end{split}                                              \tag{3.9}
\]

The point of the ideal moments is that these conclusions do not require
any estimate of \(\int f_j^4\), which could be unnecessarily large.

\hypertarget{a-globally-smooth-approximate-weyl-tensor}{%
\subsection{4. A globally smooth approximate Weyl
tensor}\label{a-globally-smooth-approximate-weyl-tensor}}

The norm of an element of \(\mathcal O\) is \(1/(2\sqrt6)\).
Consequently

\[
 \left\|q-\frac92\right\|_2^2
       \le2\|e\|_2^2+12\int_MF\le12\mathcal C.           \tag{4.1}
\]

Clipping at \(3\) moves \(q\) toward \(9/2\), so the same bound holds
with \(Q\). Since \(Q\ge3\), (3.8) implies

\[
 \|\rho-\sqrt2\|_2^2
  \le\frac8{27}S_*+\frac2{27}\left\|Q-\frac92\right\|_2^2+o_j(1).
 \tag{4.2}
\]

Here the first term follows by writing
\(\rho-\sqrt2=(2/3)(v-\sqrt Q)+(2/3)(\sqrt Q-\sqrt{9/2})\) and using
\(Q\ge3\).

To estimate direction, choose a measurable rank-one orthogonal
projection \(\Pi\) onto a lowest eigenspace of \(W^+\). Then

\[
                  A_0=I/12-\Pi/4
\]

is a nearest point of \(\mathcal O\). On the good set \(F\le1/256\),
comparison of ordered eigenvalues with \(A_0\) gives

\[
                   \lambda_2(W^+)-\lambda_1(W^+)\ge1/8.
\]

Putting \(y=(I-\Pi)\psi\), it follows from the eigenvalue error in (2.8)
that

\[
 \int_{\{F\le1/256\}}|y|^2\le\frac{16}{27}S_*+o_j(1).
\]

The complementary set has volume at most \(256\int F\), while
\(|\psi|^2\le16\). Thus

\[
                 \int_M|y|^2\le\frac{16}{27}S_*+4096\int_MF+o_j(1).
 \tag{4.3}
\]

No derivative of \(\Pi\) will be taken. Instead define the polynomial
endomorphism

\[
       A_\psi=\frac{\rho^2}{24}I-\frac18\psi\otimes\psi,
             \qquad R=W^+-A_\psi.                       \tag{4.4}
\]

It is globally smooth, trace-free, and defined even at spinor zeros. Its
eigenvalues are \(-\rho^2/12,\rho^2/24,\rho^2/24\), and
\(|A_\psi|^2=\rho^4/96\). The rank-one identity

\[
 |\psi\otimes\psi-\rho^2\Pi|^2=2\rho^2|y|^2
\]

gives the useful comparison

\[
 |A_\psi-A_0|^2
       \le\frac{\rho^2}{16}|y|^2+\frac{(\rho^2-2)^2}{48}.
 \tag{4.5}
\]

For example, (4.5) follows by inserting the intermediate matrix
\(\rho^2 A_0/2\) and applying \(|X+Y|^2\le2|X|^2+2|Y|^2\). Since
\(\rho\le4\) and \((4+\sqrt2)^2<36\), (4.1)--(4.5) imply

\[
 \begin{split}
 \|R\|_2^2
 &\le\frac{44}{27}S_*+8194\int_MF
                     +\frac19\left\|Q-\frac92\right\|_2^2+o_j(1)\\
 &\le2\Delta_++8500\mathcal C+o_j(1).
 \end{split}                                              \tag{4.6}
\]

In the final line we used \(\int e^2\le4\int|E|^2\). In particular, each
of

\[
 \|e\|_2^2,\quad\|\rho-\sqrt2\|_2^2,\quad\|R\|_2^2,
            \quad\|d\rho\|_2^2,\quad\|d\psi\|_2^2
\]

is bounded by \(C\kappa V+o_j(1)\).

\hypertarget{a-hessian-estimate-that-needs-no-sobolev-constant}{%
\subsection{5. A Hessian estimate that needs no Sobolev
constant}\label{a-hessian-estimate-that-needs-no-sobolev-constant}}

\textbf{Lemma 5.1.} Let \(u\) be a smooth section of a metric or
Hermitian vector bundle with compatible connection over a closed
Riemannian four-manifold. Let \(\mathscr F\) be its curvature, with
component norm \(\sum_{i,j}|\mathscr F_{ij}|_{\mathrm{HS}}^2\), and set
\(K=\|u\|_\infty\). Then

\[
 \begin{split}
 \|\nabla u\|_4^2&\le4K\|\nabla^2u\|_2,\\
 \|\nabla^2u\|_2^2
 &\le2\|\nabla^*\nabla u\|_2^2
       +32K^2\|\operatorname{Ric}\|_2^2
       +33K^2\|\mathscr F\|_2^2.
 \end{split}                                              \tag{5.1}
\]

\textbf{Proof.} Take real parts of Hermitian inner products when
necessary. Write \(v_i=\nabla_i u\), \(H_{ij}=\nabla_i v_j\), and
\(L=\sum_iH_{ii}=-\nabla^*\nabla u\). Integrating the divergence of
\(X_i=|v|^2\langle u,v_i\rangle\) gives

\[
 \int|v|^4
 \le K\int |v|^2(2|H|+|L|)
 \le K\|v\|_4^2(2\|H\|_2+\|L\|_2).
\]

Because \(|L|\le2|H|\) in dimension four, this proves the first
inequality, including the case \(\|v\|_4=0\).

For the second, antisymmetry gives

\[
 H_{ij}-H_{ji}=\mathscr F_{ij}u,\qquad
 \|H\|_2^2=\int\langle H_{ij},H_{ji}\rangle
                         +\frac12\sum_{i,j}\|\mathscr F_{ij}u\|_2^2.
 \tag{5.2}
\]

The commutator on the bundle-valued one-form \(v\) yields

\[
 \sum_i\nabla_i\nabla_jv_i
     =\nabla_jL+\operatorname{Ric}_{jk}v_k
                                      +\sum_i\mathscr F_{ij}v_i.
\]

Integrating by parts in the first term of (5.2) and using the
skew-adjointness of \(\mathscr F_{ij}\), we obtain the exact identity

\[
 \begin{split}
 \|H\|_2^2={}&\|L\|_2^2
  -\int\operatorname{Ric}_{ij}\langle v_i,v_j\rangle
  +\sum_{i,j}\int\langle\mathscr F_{ij}v_j,v_i\rangle
  +\frac12\sum_{i,j}\|\mathscr F_{ij}u\|_2^2.
 \end{split}                                              \tag{5.3}
\]

In particular no derivative of curvature appears. Hölder and the first
part bound the middle terms in absolute value by

\[
   4K(\|\operatorname{Ric}\|_2+\|\mathscr F\|_2)\|H\|_2.
\]

Use Young's inequality to absorb half of \(\|H\|_2^2\), and then
\((x+y)^2\le2(x^2+y^2)\). This gives exactly the second inequality in
(5.1). No Sobolev inequality or curvature smallness has been used.
\(\square\)

\hypertarget{applying-the-lemma-to-the-bounded-spinor}{%
\subsubsection{5.2. Applying the lemma to the bounded
spinor}\label{applying-the-lemma-to-the-bounded-spinor}}

The calibration budget gives

\[
 \begin{split}
 \int|W^+|^2&\le2\int F+V/12,\\
 \int|\operatorname{Ric}|^2&\le2\int|E|^2+8V,\\
 \int s^2&\le8\int|E|^2+32V.
 \end{split}                                              \tag{5.4}
\]

Equations (3.3) and (4.1) also imply

\[
                         \int f_j^2\le C(V+\mathcal C)+o_j(1).
 \tag{5.5}
\]

Indeed \(Q\le9/2+|Q-9/2|\); Cauchy--Schwarz and
\(\sqrt{V\mathcal C}\le(V+\mathcal C)/2\) suffice.

For the determinant connection, Chern--Weil and (3.2) give

\[
 \|F_A\|_2^2
   =2\|F_A^+\|_2^2-4\pi^2a^2
   \le4\int f_j^2+4\pi^2(-a^2)_+.
 \tag{5.6}
\]

The positive-spinor connection curvature \(\mathscr F_A\) is a fixed
linear representation of the Riemannian and determinant curvatures. In
the block decomposition of the Riemannian curvature operator, its
nontrivial spin action uses the \(\Lambda^+\) output: it sees \(W^+\),
\(s\), and the trace-free Ricci block, but not \(W^-\). Consequently
(5.4)--(5.6) imply

\[
             \int|\mathscr F_A|^2
               \le C\bigl(V+\mathcal C+(-a^2)_+\bigr)+o_j(1).
 \tag{5.7}
\]

From (2.4) and \(\rho\le4\),

\[
 \|\nabla_A^*\nabla_A\Phi\|_2^2
       \le\frac12\int s^2+8\int f_j^2
       \le C(V+\mathcal C)+o_j(1).                        \tag{5.8}
\]

Apply Lemma 5.1 with \(\|\Phi\|_\infty\le2\). Since \(A_\psi\) is a
parallel polynomial of degree four in \(\Phi\),

\[
 |\nabla^2A_\psi|
      \le C\bigl(|\Phi|^3|\nabla_A^2\Phi|
                               +|\Phi|^2|\nabla_A\Phi|^2\bigr).
\]

We conclude that

\[
 \int\bigl(|\nabla_A^2\Phi|^2+|\nabla_A\Phi|^4
                                      +|\nabla^2A_\psi|^2\bigr)
       \le C\bigl(V+\mathcal C+(-a^2)_+\bigr)+o_j(1).
 \tag{5.9}
\]

This proves the required second-derivative budget without
differentiating the weight and without an estimate of the anti-self-dual
Weyl curvature.

\hypertarget{a-weak-bianchi-self-test-and-its-incompatible-lower-bound}{%
\subsection{6. A weak Bianchi self-test and its incompatible lower
bound}\label{a-weak-bianchi-self-test-and-its-incompatible-lower-bound}}

\hypertarget{tensor-normalization-and-the-bianchi-upper-bound}{%
\subsubsection{6.1. Tensor normalization and the Bianchi upper
bound}\label{tensor-normalization-and-the-bianchi-upper-bound}}

For an orthonormal basis \(\epsilon_a\) of \(\Lambda^+\), put
\(\omega_a=\sqrt2\epsilon_a\). Associate to a symmetric endomorphism
\(A\) the four-tensor

\[
 W(A)_{ijkl}=\frac12\sum_{a,b}A_{ab}(\omega_a)_{ij}(\omega_b)_{kl}.
 \tag{6.1}
\]

Let

\[
                    J_\psi=\operatorname{div}W(A_\psi),
\]

using the full three-index norm \(\sum_{j,k,l}|J_{jkl}|^2\). Changing
the overall sign of divergence changes none of the following squared
norms.

The contracted Bianchi identity expresses \(\operatorname{div}W^+\) as a
parallel algebraic contraction of \(\nabla\operatorname{Ric}^\circ\) and
\(\nabla s\). This follows by contracting the Riemannian Bianchi
identity, subtracting the Schouten part, and projecting the final
two-form pair onto \(\Lambda^+\); that projection is parallel. Since
\(\nabla\operatorname{Ric}=\nabla E\) and
\(\nabla s=\nabla\operatorname{tr}E\), integration by parts and
\(W(A_\psi)=W^+-W(R)\) show, for every smooth test tensor \(Z\), that

\[
 \left|\int\langle J_\psi,Z\rangle\right|
       \le C(\|E\|_2+\|R\|_2)\|\nabla Z\|_2.
 \tag{6.2}
\]

Set \(Z=J_\psi\). It is smooth, and
\(\|\nabla J_\psi\|_2\le C\|\nabla^2A_\psi\|_2\). Equations (4.6) and
(5.9) therefore give

\[
 \frac1V\int|J_\psi|^2
            \le C\sqrt\kappa\sqrt{1+\kappa+B_a}+o_j(1).
 \tag{6.3}
\]

This is a weak identity tested against a smooth tensor already supplied
by the spinor. It asserts no pointwise bound on
\(\nabla\operatorname{Ric}\).

\hypertarget{the-exact-divergence-identity}{%
\subsubsection{6.2. The exact divergence
identity}\label{the-exact-divergence-identity}}

Define

\[
 B_\psi=\frac{s}{3}\rho^2-2\langle W^+\psi,\psi\rangle,
 \qquad C_\psi=\langle\iota_{\nabla\rho}\psi,*d\psi\rangle.
\]

We now derive the identity used to lower-bound the left side of (6.3).
Pairing the self-dual Hodge--Weitzenböck equation with \(\rho^2\psi\)
and integrating by parts gives

\[
 \int\rho^2\bigl(|\nabla\psi|^2+B_\psi-2|d\psi|^2\bigr)
          =4\int\rho C_\psi-2\int\rho^2|d\rho|^2.
 \tag{6.4}
\]

To see the mixed terms explicitly, use \(d^*\psi=-*d\psi\),
\(d^*(\rho^2\psi)=\rho^2d^*\psi-2\rho\iota_{\nabla\rho}\psi\), and
\(d(\rho^2\psi)=\rho^2d\psi+2\rho d\rho\wedge\psi\). The two
exterior-derivative mixed terms add to \(4\rho C_\psi\); the
rough-Laplacian mixed term is \(2\rho^2|d\rho|^2\).

The required pointwise first-jet identities are

\[
 |\nabla A_\psi|^2
       =\frac{\rho^2}{32}|\nabla\psi|^2
                                +\frac{\rho^2}{96}|d\rho|^2,
 \tag{6.5}
\]

\[
 |\nabla A_\psi|^2-|J_\psi|^2
    =\frac{\rho^2}{64}(|\nabla\psi|^2-2|d\psi|^2)
      +\frac{\rho^2}{32}|d\rho|^2-\frac\rho{16}C_\psi.
 \tag{6.6}
\]

Here is a direct component verification that fixes all normalization
factors. At a point with \(\rho>0\), choose a parallel orthonormal frame
so that

\[
 \begin{split}
 \epsilon_1&=(dx^0\wedge dx^1+dx^2\wedge dx^3)/\sqrt2,\\
 \epsilon_2&=(dx^0\wedge dx^2+dx^3\wedge dx^1)/\sqrt2,\\
 \epsilon_3&=(dx^0\wedge dx^3+dx^1\wedge dx^2)/\sqrt2,\\
 \psi&=\rho\epsilon_1,\qquad
 \nabla_m\psi=c_m\epsilon_1+a_m\epsilon_2+b_m\epsilon_3.
 \end{split}
\]

Then \(c_m=\partial_m\rho\), and differentiating (4.4) gives

\[
 \nabla_m A_\psi=
 \begin{pmatrix}
  -\rho c_m/6&-\rho a_m/8&-\rho b_m/8\\
  -\rho a_m/8&\rho c_m/12&0\\
  -\rho b_m/8&0&\rho c_m/12
 \end{pmatrix}.                                         \tag{6.7}
\]

The squared matrix norm gives (6.5). For (6.6), insert (6.7) in
\(J_{jkl}=\sum_{m,a,b}(\nabla_mA)_{ab}(\epsilon_a)_{mj}(\epsilon_b)_{kl}\),
use \(\sum_{k,l}(\epsilon_b)_{kl}(\epsilon_c)_{kl}=2\delta_{bc}\), and
compute
\(d\psi_{ijk}=\nabla_i\psi_{jk}+\nabla_j\psi_{ki} +\nabla_k\psi_{ij}\)
in the displayed basis. Explicitly,

\[
 *d\psi=\frac1{\sqrt2}
 \begin{pmatrix}
  -a_2-b_3-c_1\\
   a_3-b_2+c_0\\
   a_0+b_1-c_3\\
  -a_1+b_0+c_2
 \end{pmatrix},
 \qquad
 \iota_{\nabla\rho}\psi
       =\frac{\rho}{\sqrt2}(-c_1,c_0,-c_3,c_2).
\]

These expressions and (6.7) give (6.6) by summing squares and collecting
the mixed terms in \(C_\psi\). Thus (6.5)--(6.6) are algebraic
identities in the twelve arbitrary numbers \(c_m,a_m,b_m\), with no
field equation imposed. At a spinor zero, both \(\psi\) and its first
derivative vanish, and \(d\rho=0\); the identities extend there as well.

Combining (6.4)--(6.6) eliminates \(\int\rho^2|\nabla\psi|^2\) and gives

\[
 \begin{split}
 \int|J_\psi|^2={}&-\frac1{64}\int\rho^2 B_\psi
       +\frac1{16}\int\rho^2|d\psi|^2
       -\frac5{96}\int\rho^2|d\rho|^2
       +\frac18\int\rho C_\psi.
 \end{split}                                              \tag{6.8}
\]

\hypertarget{the-positive-lower-bound-and-contradiction}{%
\subsubsection{6.3. The positive lower bound and
contradiction}\label{the-positive-lower-bound-and-contradiction}}

Because \(A_\psi\psi=-\rho^2\psi/12\), the curvature term in (6.8) is

\[
 -\frac{\rho^2B_\psi}{64}
        =-\frac{s\rho^4}{192}-\frac{\rho^6}{384}
                                +\frac{\rho^2\langle R\psi,\psi\rangle}{32}.
 \tag{6.9}
\]

At \(s=-4\), \(\rho=\sqrt2\), and \(R=0\), it equals \(1/16\). On the
interval \(0\le\rho\le4\), the error in (6.9) is bounded by

\[
                  C\bigl(|e|+|\rho-\sqrt2|+|R|\bigr).
\]

Also \(|C_\psi|\le\rho|d\rho||d\psi|\). Discard the nonnegative
\(|d\psi|^2\) term in (6.8), apply (3.9), (4.2), and (4.6), and divide
by \(V\). Cauchy--Schwarz gives

\[
 \frac1V\int|J_\psi|^2
                   \ge\frac1{16}-C(\sqrt\kappa+\kappa)+o_j(1).
 \tag{6.10}
\]

For each fixed metric let the smooth-weight approximation errors tend to
zero in (6.3) and (6.10). The two numerical inequalities remain valid
without choosing any limit of the spinors. For a sufficiently small
absolute \(\kappa_0\le1\), the lower bound is at least \(1/32\) when
\(\kappa\le\kappa_0\). The upper bound then forces

\[
                         \kappa(1+\kappa+B_a)\ge c_1>0.
\]

As \(1+\kappa+B_a\le2(1+B_a)\) in this range, (1.8) follows. For
\(\kappa>\kappa_0\), it follows after decreasing the same universal
constant. This proves Theorem 1.3. \(\square\)

\hypertarget{from-the-calibration-gap-to-topological-integral-rigidity}{%
\subsection{7. From the calibration gap to topological integral
rigidity}\label{from-the-calibration-gap-to-topological-integral-rigidity}}

We now prove Theorem 1.2, first at the negative normalization. Assume
(1.5), and put

\[
 r=\frac1V\int|\operatorname{Ric}+g|^2,\qquad
 d=\left(\frac TA-\frac23\right)_+.
\]

Continue to use \(\alpha=\sqrt{2/3}|W^+|\) and
\(w=\lambda_{\min}(W^+)\), and introduce

\[
 t=s+24\alpha=s+8\sqrt6|W^+|,\qquad
 b=2s/3+2w,\qquad
 G=(2s/3-2\alpha)^2-b_-^2.                               \tag{7.1}
\]

Since \(b\ge2s/3-2\alpha\), we have \(G\ge0\), including where the
scalar curvature is positive. Direct expansion gives

\[
 t^2+18G=216\left(s^2/24+2|W^+|^2\right)-18b_-^2.
\]

Apply (1.1) and (2.12) to obtain

\[
                    \int(t^2+18G)\le108P+864\pi^2Ad.
 \tag{7.2}
\]

Here \(p\ge A\) and \(T-2p/3\le T-2A/3\le Ad\). Since
\(\int|E|^2=P+\frac14\int e^2\), (7.2) implies

\[
             \int(|E|^2+e^2+t^2+G)\le109rV+864\pi^2Ad.
 \tag{7.3}
\]

We record the pointwise orbit comparison, including its scalar-positive
branch. Set \(z=w+\alpha\); trace-freeness gives \(0\le z\le\alpha\),
and minimizing over the orbit gives exactly

\[
 F=\frac32(\alpha-1/6)^2+z/2,\qquad
                         \alpha=(t-e+4)/24.              \tag{7.4}
\]

If \(s\le-3\), then \(b<0\) and \(G=z(-8s/3+8\alpha-4z)\ge8z\). Thus

\[
                       F\le(t^2+e^2)/192+G/16.
\]

If \(s>-3\), then \(e>1\). Use \(z\le\alpha\), \(4-e\le3e^2\), and
\(t\le(t^2+1)/2\) in (7.4) to get

\[
                       F\le(t^2+5e^2)/64.
\]

In both cases,

\[
                             F\le e^2+t^2+G.             \tag{7.5}
\]

The second completed square is

\[
 27\left(s^2/24+2|W^+|^2\right)-q^2=t^2/8.
\]

It and (1.1) yield

\[
 \Delta_+\le\frac{27}{2}rV+108\pi^2Ad.
\]

Combining this with (7.3)--(7.5),

\[
                     \kappa\le\frac{245}{2}r
                                      +972\pi^2\frac AVd.
 \tag{7.6}
\]

The scalar estimate (2.13) gives

\[
 32\pi^2A\le\int s^2\le32V+8rV,
              \qquad \frac AV\le\frac{4+r}{4\pi^2}.
 \tag{7.7}
\]

For \(r\le1\), therefore,

\[
                              \kappa\le1215(r+d).        \tag{7.8}
\]

When \(r\le1/4\), reverse triangle inequality in \(L^2\) and (1.1) also
imply

\[
 4\pi^2T
     \ge\left[\frac{(4-2\sqrt r)^2}{24}-\frac r2\right]V
     \ge V/4>0.                                          \tag{7.9}
\]

The square condition \(a^2\ge T\) then gives \(B_a=0\). Apply Theorem
1.3 and (7.8). If \(r>1/4\), no argument is needed. In either case

\[
       r+d\ge\varepsilon_0,
       \qquad \varepsilon_0:=\min\{1/4,c_{\mathrm{cal}}/1215\}.
 \tag{7.10}
\]

For reference, a weaker square hypothesis has a correspondingly
nonuniversal version. If \(a^2\ge-C_AA\), \(C_A\ge0\), the same proof
gives

\[
 r+d\ge\min\left\{1,
          \frac{c_{\mathrm{cal}}}{1215(1+5C_A/(4\pi^2))}\right\}.
 \tag{7.11}
\]

If instead \(a^2\ge T-C_AA\), replace the first entry of that minimum by
\(1/4\), using (7.9). An arbitrary monopole class must not silently be
assumed to have \(a^2\ge T\).

\hypertarget{positive-scale-arbitrary-scale-and-the-ricci-flat-case}{%
\subsection{8. Positive scale, arbitrary scale, and the Ricci-flat
case}\label{positive-scale-arbitrary-scale-and-the-ricci-flat-case}}

\hypertarget{the-positive-normalization}{%
\subsubsection{8.1. The positive
normalization}\label{the-positive-normalization}}

Put \(r_+=V^{-1}\int|\operatorname{Ric}-g|^2\). Since
\(s_-^2\le(s-4)^2\le4|\operatorname{Ric}-g|^2\), (2.13) gives

\[
                            A/V\le r_+/(8\pi^2).
\]

If \(r_+\le1/4\), the analogue of (7.9) gives \(4\pi^2T/V\ge1/4\). If
also \(d\le1/4\), then \(T/A\le2/3+d\), so

\[
  \frac14\le\frac{4\pi^2T}{V}
             \le\frac{r_+}{2}\left(\frac23+d\right)
             \le\frac{11}{96},
\]

a contradiction. In particular \(r_++d>1/4\).

\hypertarget{scale-optimization}{%
\subsubsection{8.2. Scale optimization}\label{scale-optimization}}

For \(\lambda\ne0\), apply the appropriate normalized result to
\(|\lambda|g\). The positive harmonic subspace and the conformal
solvability hypothesis are unchanged. The normalized error becomes

\[
 \frac{\int|\operatorname{Ric}_g-\lambda g|^2}{\lambda^2V}.
\]

Taking the infimum proves (1.6). This completes Theorem 1.2. \(\square\)

For interpretation, write

\[
 R_2=\int|\operatorname{Ric}|^2,\qquad S=\int s,\qquad
 L_s=\int(s-S/V)^2.
\]

When \(R_2>0\), direct minimization in \(1/\lambda\) gives

\[
              \mathcal E(g)=4-\frac{S^2}{VR_2}
                             =\frac{4P+L_s}{R_2}.        \tag{8.1}
\]

This formula includes \(S=0\), when the infimum is \(4\) and need not be
attained. If \(R_2=0\), definition (1.3) gives \(\mathcal E(g)=4\), not
zero. Thus a lower bound for \(\mathcal E\) alone does not exclude a
Ricci-flat metric. Ricci-flat metrics in Theorem 1.2 are excluded
separately by (2.13), since \(p\ge A>0\).

For an Einstein metric with nonzero Einstein constant, \(\mathcal E=0\).
A pointwise operator error also dominates the integral one: for
\(\sigma=\pm1\),

\[
 \delta_\sigma(g):=\|\operatorname{Ric}^{\sharp}-\sigma I\|_{L^\infty,\mathrm{op}}
      \quad\Longrightarrow\quad\mathcal E(g)\le4\delta_\sigma(g)^2.
 \tag{8.2}
\]

There are exactly four Ricci eigenvalues, which explains this factor.

\hypertarget{metricwise-polarized-classes-and-chamber-bookkeeping}{%
\subsection{9. Metricwise polarized classes and chamber
bookkeeping}\label{metricwise-polarized-classes-and-chamber-bookkeeping}}

The analytic theorem needs the quantifiers

\[
 \boxed{\begin{gathered}
   \text{for every }g\text{ there exists a Spin}^c
       \text{ structure }\mathfrak s_g\\
   \text{such that for every }f>0,\text{ equation (1.2)}\\
   \text{has a solution for the same }\mathfrak s_g.
 \end{gathered}}
 \tag{9.1}
\]

It does not require one universal choice of signs for all metrics. This
distinction is essential when \(b_+=1\).

\hypertarget{the-connected-sum-input-with-insertions}{%
\subsubsection{9.1. The connected-sum input with
insertions}\label{the-connected-sum-input-with-insertions}}

Let \(X\) have \(b_+(X)>0\), and suppose a Spin\(^c\) structure
\(\mathfrak s\), with real first Chern class \(c\), has a nonzero
Seiberg--Witten invariant evaluated on a homogeneous insertion \(z\).
The invariant is defined on

\[
 \mathbb A(X)=\mathbb Z[U]\otimes\Lambda^*H_1(X;\mathbb Z)/\mathrm{torsion},
                 \qquad\deg U=2,\quad\deg H_1=1.
\]

Its expected dimension is \((c^2-T_X)/4\), where
\(T_X=2\chi(X)+3\tau(X)\). Assume

\[
                    T_X>0,\qquad A:=c^2\ge T_X.          \tag{9.2}
\]

The second inequality follows automatically from the nonnegative degree
of a nonzero ordinary invariant, but is displayed because it is used in
the proof. For \(b_+=1\), the nonzero invariant here must be in the
small-perturbation chamber at the relevant polarization, not merely in
some unspecified chamber.

Let

\[
 M=X\#k\,\overline{\mathbb{CP}}^{2}\#\ell(S^1\times S^3),
 \qquad D=k+4\ell,\qquad T_M=T_X-D.                      \tag{9.3}
\]

For a prescribed positive polarization \(H=H_g^+\), choose the signs of
the exceptional generators \(E_j\) so that

\[
                         c_H^+\mathbin\cdot E_j\le0,
\]

and set

\[
                             a_H=c-\sum_{j=1}^kE_j.
 \tag{9.4}
\]

Then

\[
 \begin{split}
 (a_H^+)^2
   &=(c_H^+)^2-2\sum_jc_H^+\mathbin\cdot E_j
                              +\left|\sum_j(E_j)_H^+\right|^2
      \ge(c_H^+)^2\ge c^2=A,\\
 a_H^2&=A-k\ge T_X-k-4\ell=T_M.
 \end{split}                                              \tag{9.5}
\]

The ordinary blow-up formula with exceptional coefficient \(\pm1\)
preserves expected dimension and preserves the original insertion \(z\).
The \(S^1\times S^3\) operation raises expected dimension by one and
adds the degree-one insertion of its new circle. After \(\ell\)
operations the nonzero invariant is therefore evaluated on

\[
                           z\,\mu(C_1)\cdots\mu(C_\ell).
 \tag{9.6}
\]

These statements include positive-dimensional base invariants. They are
the special cases of the blow-up and surgery formulas in
\href{https://arxiv.org/pdf/math/0005268}{Ozsváth--Szabó, Theorem 3.2,
Proposition 2.2 and Remark 2.3}. The last remark supplies the
chamber-compatible version when \(b_+=1\). Using only a zero-dimensional
numerical invariant and omitting the insertions in (9.6) would not
justify the general statement.

\hypertarget{what-the-chamber-condition-means}{%
\subsubsection{9.2. What the chamber condition
means}\label{what-the-chamber-condition-means}}

For \(b_+(X)=1\), choose a unit forward generator \(\omega_H\) of the
positive line on \(M\) so that \(c\cdot\omega_H>0\). The \(X\) component
\(\omega_X\) has positive square, lies in the corresponding forward
cone, and satisfies \(c\cdot\omega_X>0\), because \(c^2>0\). The chosen
exceptional signs give \(E_j\cdot\omega_H\le0\), hence

\[
                          a_H\cdot\omega_H>0.             \tag{9.7}
\]

Thus the small-perturbation discriminants for the base and the selected
blow-up structure have the same sign in corresponding chamber
conventions. Adding \(S^1\times S^3\) leaves the second-cohomology
polarization unchanged. The gluing formulas consequently give a nonzero
invariant in the chamber containing zero perturbation for this selected
class, provided that was the specified nonzero base chamber.

For any metric with this polarization, take generic perturbations
tending to zero in that chamber. Nonzero invariant gives solutions, and
the usual fixed-metric Seiberg--Witten compactness gives an unperturbed
solution. A reducible limit would have \(F_A^+=0\), contradicting
\((a_H^+)^2\ge A>0\). Every metric conformal to \(g\) has the same
polarization and the same nonzero discriminant sign. Thus the
\emph{same} class (9.4) satisfies Definition 1.1 for every weight, as
required by (9.1).

For \(b_+>1\), there is no chamber ambiguity in this argument. For
\(b_+=1\), the statement must not be replaced by the assertion that
every exceptional sign choice is an ordinary all-metric monopole class.

\hypertarget{the-minimal-general-type-case}{%
\subsubsection{9.3. The minimal general-type
case}\label{the-minimal-general-type-case}}

If \(X\) is a minimal compact complex surface of general type with its
complex orientation, then

\[
                          A=c_1(X)^2=T_X>0.
\]

The canonical Spin\(^c\) structure has a nonzero Seiberg--Witten
invariant. For \(b_+>1\), the standard Kähler/symplectic nonvanishing
statement supplies the input directly. For \(b_+=1\), the canonical
class \(K\) has \(K^2>0\) and positive pairing with a Kähler form. It is
therefore in the forward time-like cone and pairs positively with every
forward period ray. For \(c=-K\), the small-perturbation chamber agrees
with the canonical nonzero chamber, up to the simultaneous choice of
sign convention for the chamber discriminant. This is the
light-cone/chamber argument underlying
\href{https://arxiv.org/pdf/math/0003068}{LeBrun's polarized
construction, Proposition 3.2}. The chamber convention and canonical
nonvanishing can also be read in
\href{https://arxiv.org/pdf/math/0012048}{Li--Liu, §2.1 and Theorem
2.2}.

The positivity of \(K\) on the symplectic or Kähler form matters:
\(b_+=1\) and \(K^2>0\) alone do not identify the nonzero
small-perturbation chamber for an arbitrary symplectic manifold.

\hypertarget{the-einstein-obstruction-and-its-precise-relation-to-lebrun}{%
\subsection{10. The Einstein obstruction and its precise relation to
LeBrun}\label{the-einstein-obstruction-and-its-precise-relation-to-lebrun}}

\textbf{Theorem 10.1 (connected-sum integral obstruction).} Under the
gauge-theoretic and chamber assumptions of §9.1, every metric on (9.3)
satisfies

\[
 \mathcal E(g)+\left(\frac{T_M}{A}-\frac23\right)_+
                         \ge\varepsilon_0.               \tag{10.1}
\]

It also satisfies the slightly weaker but base-topological statement

\[
 \boxed{\displaystyle
 \mathcal E(g)+\left(\frac13-\frac{D}{T_X}\right)_+
                         \ge\varepsilon_0.}              \tag{10.2}
\]

For minimal general type, \(A=T_X=c_1(X)^2\), so the two statements
coincide.

\textbf{Proof.} Equation (9.5) and the conformal solvability just proved
give all the hypotheses of Theorem 1.2. If \(T_M>0\), then \(A\ge T_X\)
implies \(T_M/A\le T_M/T_X\), yielding (10.2). If \(T_M\le0\), both
positive-part terms are zero. This case split avoids reversing an
inequality by dividing a negative numerator by the larger denominator.
\(\square\)

In particular, put \(\gamma=\varepsilon_0/2\). Then

\[
             D\ge(1/3-\gamma)T_X
 \quad\Longrightarrow\quad
             \mathcal E(g)\ge\varepsilon_0/2
             \quad\hbox{for every metric }g.             \tag{10.3}
\]

There is no Einstein metric in this band: nonzero Einstein constants
contradict (10.3), and zero Einstein constant contradicts (2.13). At
either normalized sign, (8.2) also gives

\[
                    \delta_\sigma(g)\ge\sqrt{\varepsilon_0/8}.
 \tag{10.4}
\]

An Einstein metric would in fact require \(D/T_X\le1/3-\varepsilon_0\).
The equality endpoint of that necessary condition is not excluded by
(10.2) alone; the closed band (10.3) avoids that endpoint issue.

\hypertarget{exact-source-scope}{%
\subsubsection{10.2. Exact source scope}\label{exact-source-scope}}

The version of \href{https://arxiv.org/pdf/math/0003068}{LeBrun's paper
verified here} is arXiv:math/0003068v1, March 2000. Its Proposition 3.2
constructs the metricwise polarized classes on the connected sums in
(9.3). Its Theorem 3.3 states the classical threshold

\[
                              k+4\ell\ge T_X/3,
\]

and Corollary 3.4 specializes the result to blow-ups of minimal complex
surfaces of general type. We use the ordinary nontrivial Seiberg--Witten
input, including its insertion and chamber meaning, not merely the
topology of a manifold with the same Euler characteristic and signature.
The proof above handles the positive-budget case of the general theorem,
not only its general-type corollary.

The numerical band below the classical threshold is nonvacuous in
families with unbounded budget. For example, let
\(X_m=\Sigma_2\times\Sigma_m\), \(m\ge2\). Then \(A=T_{X_m}=8(m-1)\).
With \(\ell=0\) and \(k=\lceil A/3\rceil-1\),

\[
                       0<1/3-k/A\le1/A\longrightarrow0.
\]

For sufficiently large \(m\), these ratios lie in (10.3). Since the
universal analytic constant has not been made numerical, this argument
does not specify the first qualifying \(m\).

\hypertarget{nonpositive-budgets-and-the-exceptional-zero-corner}{%
\subsubsection{10.3. Nonpositive budgets and the exceptional zero
corner}\label{nonpositive-budgets-and-the-exceptional-zero-corner}}

There is a direct all-metric formula, independent of monopoles. For
\(R_2>0\), (1.1) and (8.1) give

\[
 \mathcal E(g)
   =1+\frac{12\int|W^+|^2-24\pi^2T+L_s}{R_2}.           \tag{10.5}
\]

Thus \(T\le0\) implies \(\mathcal E(g)\ge1\). For a Ricci-flat metric
its value is \(4\) by definition, so the same lower bound holds. This
lower bound is not a Ricci-flat exclusion.

For an Einstein metric, (1.1) reduces to

\[
               4\pi^2T=\int(s^2/24+2|W^+|^2)\ge0.
\]

Consequently, when \(T_X<0\), or when \(T_X=0\) and \(D>0\), the
connected sum has \(T_M<0\) and admits no Einstein metric. These are
classical Chern--Gauss--Bonnet consequences and require no
positive-budget monopole argument.

When \(T_X=D=0\), however, equality permits Ricci-flat, anti-self-dual
metrics in the chosen orientation. K3 and the four-torus provide the
standard zero-budget examples. The unqualified statement of Theorem 3.3
in the verified preprint includes this corner, but its appeal to the
Hitchin--Thorpe inequality does not exclude these equality cases. We
therefore do \textbf{not} claim an unqualified subsumption of that
literal zero-corner wording, nor infer an exception-free statement for a
separately unverified published version. The conclusions established
here are: the positive-budget integral strengthening with the stated
gauge input; the classical negative-\(T_M\) exclusion; and an explicit
reservation of the Ricci-flat zero corner.

\hypertarget{logical-dependencies-and-limitations}{%
\subsection{11. Logical dependencies and
limitations}\label{logical-dependencies-and-limitations}}

The analytic proof has the following complete chain:

\[
 \begin{gathered}
 \text{one actual conformal Spin}^c\text{ structure}
 \ \Longrightarrow\ \text{solutions for smooth truncated weights}\\
 \Longrightarrow\ \rho\le4\text{ and small retained first-order defects}
 \ \Longrightarrow\ W^+\approx A_\psi\text{ with }A_\psi\text{ smooth}\\
 \Longrightarrow\ \text{bundle Hessian budget and weak Bianchi upper bound}\\
 \Longrightarrow\ \text{incompatibility with the exact positive divergence term}.
 \end{gathered}
\]

The proof neither differentiates an eigenprojection nor normalizes
\(\psi\) across its zero set. Smooth approximation of weights is used
only for numerical inequalities at a fixed metric; there is no claimed
uniform regularity or limiting solution for a nonsmooth weight. The
Hessian estimate uses the positive-spinor connection, which is why no
\(L^2\) budget for \(W^-\) enters. The topological theorem additionally
needs the class-square bound and the actual class/chamber construction;
the scalar curvature inequality by itself is not a substitute for those
inputs.

This argument proves an obstruction and an integral gap. It does not
prove convergence of near-Einstein metrics, classify all equality
metrics of other curvature functionals, produce an Einstein metric, or
settle arbitrary pointwise Ricci-pinching problems without the
monopole/topological hypotheses displayed above.

\hypertarget{references-and-provenance}{%
\subsection{References and provenance}\label{references-and-provenance}}

\begin{itemize}
\tightlist
\item
  Claude LeBrun, \emph{Ricci Curvature, Minimal Volumes, and
  Seiberg--Witten Theory},
  \href{https://arxiv.org/pdf/math/0003068}{arXiv:math/0003068v1},
  especially Proposition 3.2, Theorem 3.3 and Corollary 3.4. The source
  version and the zero-budget caveat are specified in §10.
\item
  Peter Ozsváth and Zoltán Szabó, \emph{Higher Type Adjunction
  Inequalities in Seiberg--Witten Theory},
  \href{https://arxiv.org/pdf/math/0005268}{arXiv:math/0005268},
  especially Proposition 2.2, Remark 2.3 and Theorem 3.2. These supply
  insertion-aware surgery and blow-up formulas.
\item
  Tian-Jun Li and Ai-Ko Liu, \emph{Uniqueness of Symplectic Canonical
  Class, Surface Cone and Symplectic Cone of 4-Manifolds with
  \(b^+=1\)},
  \href{https://arxiv.org/pdf/math/0012048}{arXiv:math/0012048}, §2.1
  and Theorem 2.2, for chamber conventions and canonical nonvanishing.
\item
  Claude LeBrun, \emph{Four-Manifolds, Curvature Bounds, and Convex
  Geometry}, \href{https://arxiv.org/pdf/math/0611450}{arXiv preprint},
  especially Lemma 3.7, for the conformal weighted-equation input.
\end{itemize}

This text consolidates the local weighted-monopole,
integral-calibration, all-sign, chamber-selection and LeBrun-comparison
arguments. Those earlier working records are provenance, not proof
dependencies of the analytic argument presented here. Numerical and
symbolic checks are useful for detecting transcription errors; none is
used in place of a mathematical identity or a gauge-theoretic existence
theorem.

\end{document}
